\documentclass[french]{book}
\usepackage{teach}
\usepackage{verbatim}
\usepackage{amsmath,amsfonts,amssymb}
\usepackage{pgfplots}
%\usepackage{geometry}
%\geometry{top=1cm, bottom=1cm, left=2cm, right=2cm}
\usepackage[utf8]{inputenc}
\usepackage[french]{babel}
\redefinecolor{thm}{title}{bgcolor}{alizarin}
\redefinecolor{prop}{title}{bgcolor}{meatbrown}
\redefinecolor{defn}{title}{bgcolor}{olivine}
\definecolor{alizarin}{rgb}{0.82, 0.1, 0.26}
\definecolor{meatbrown}{rgb}{0.9, 0.72, 0.23}
\definecolor{olivine}{rgb}{0.6, 0.73, 0.45}
\usepackage{pgf,tikz,tkz-tab}
\pgfplotsset{compat=newest}
\usepackage{mathrsfs}
\usetikzlibrary{arrows}

\begin{document}



\newpage


\begin{center}
\huge{ Devoir surveillé VI}
\end{center}
%\title{Devoir surveillé V}
Nom: 

Prénom: 

Date: 31/01/20

\vspace{1cm}
\pagestyle{empty}

Soit $(O,\overrightarrow{i},\overrightarrow{j})$ un repère orthonormé.

\definecolor{xdxdff}{rgb}{0.49019607843137253,0.49019607843137253,1}
\begin{tikzpicture}[scale=1,line cap=round,line join=round,>=triangle 45,x=1cm,y=1cm]
\begin{axis}[
x=1cm,y=1cm,
axis lines=middle,
ymajorgrids=true,
xmajorgrids=true,
xmin=-2.45768508267049,
xmax=14.496651298051097,
ymin=-2.7760577599007705,
ymax=11.422248603775223,
xtick={-2,-1,...,14},
ytick={-2,-1,...,11},]
\clip(-2.45768508267049,-2.7760577599007705) rectangle (14.496651298051097,11.422248603775223);
\draw [->,line width=1pt] (8,7) -- (1,3);
\draw [->,line width=1pt] (0,0) -- (1,0);
\draw [->,line width=1pt] (0,0) -- (0,1);
\begin{scriptsize}
\draw[color=black] (4.707992961647939,4.798273011083009) node {$\overrightarrow{u}$};
\draw[color=black] (0.4884159700333731,-0.23860943386228955) node {$\overrightarrow{i}$};
\draw[color=black] (-0.3669036904290388,0.5786960194684571) node {$\overrightarrow{j}$};
\draw [fill=xdxdff] (0,0) circle (2.5pt);
\draw[color=xdxdff] (-0.4619392082581957,-0.3146378481256148) node {$O$};
\end{scriptsize}
\end{axis}
\end{tikzpicture}

\section*{Application du cours}

\begin{enumerate}

\item Dans ce repère, donner les coordonnées du vecteur $\overrightarrow{u}$.

\item Représenter dans le repère les points $A(-2 ;6)$ et $B(1;2)$.

\item Déterminer les coordonnées  du vecteur $\overrightarrow{AB}$.

\item Déterminer les coordonnées du vecteur $\overrightarrow{u} + \overrightarrow{AB} $

\item Soit $k \in \mathbb{R}$, déterminer, en fonction de $k$, les coordonnées du vecteur $k\overrightarrow{AB}$. 

\item On pose $C(0;3)$ et $D(2,0)$. Déterminer la longueur du segment $[CD]$.

\item Soit $M$ le milieu du segment $[CD]$, déterminer les coordonnées du point $M$.   

\item Soit $\overrightarrow{EF} \left( \begin{array}{c}
-4 \\
6 
\end{array} \right)$
 Montrer que les droites $(EF)$ et $(DC)$ sont parallèles.

\end{enumerate}

\section*{Problème}

\textit{(Dans ce problème, il est conseillé de faire un dessin au brouillon)}\\

Soit $L$ un point du plan. On note $\mathcal{C}$ le cercle de centre $L$ et de rayon $r$, où $r\in\mathbb{R}_+^*$.\\

On rappel que $\mathcal{C}$ est l’ensemble des points qui sont à une distance $r$ du point $L$. \\

On pose $L(3;5)$ et $r=3$. On souhaite qu’un point $M$ appartiennent au $\mathcal{C}$. 
\begin{enumerate}
\item Recopier et compléter l’assertion suivante :

$$ M \in \mathcal{C} \Longleftrightarrow \vert \vert  \cdots \vert \vert =  .... $$

On pose M$(x;y)$ avec $x,y \in \mathbb{R}$. 


\item Déterminer les coordonnées du vecteur $\overrightarrow{LM}$.
\item En déduire alors une équation qui traduit le fait que $M \in \mathcal{C}$.
\end{enumerate}







\section*{Calcul littéral}

Développer et réduire chacune des expressions.\\

\begin{enumerate}
\item[] $A=(3x-2)(6x-3)$
\item[]$B=(5-4x)^2$
\item[] $C=(3x+5)^2$
\item[] $D=(2x+1)(2x-1)$\\
\end{enumerate} 

Déterminer pour quel(s) valeur(s) de on a $x$, $A=0$ et $B=0$.\\


Vrai ou faux? Justifier

$$ (\sqrt{3}-\sqrt{2})(\sqrt{3}+\sqrt{2})=1$$\\

Déterminer pour quel(s) valeur(s) de on a $x$, $C=2$
\end{document}